%%%%%%%%%%%%%%%%%%%%%%%%%%%%%%%%%%%%%%%%%%%%%%%%%%%%
\section{Experiments}

All experiments were conducted on a single machine
(Apple M3 Pro, 18 GB RAM) using Rust~1.82 release builds.
The CASAS dataset~\cite{casas2013} records timestamped ON/OFF events
from motion and door sensors in a smart home,
discretized into 1-hour slots.

\textbf{Benchmark methodology.}
In production, computation time and transaction arrival time are
interleaved: the system waits for $n$ transactions, then computes.
To measure \textit{pure computation time} independently of arrival
latency, we pre-load the full dataset and feed it to the system
$n$ transactions at a time, triggering each scan immediately after
the previous one completes.
This isolates algorithmic cost from network or sensor delays and
yields a fair comparison across different values of $n$.

\subsection{EX-1: Scan Frequency vs.\ Processing Speed}

We varied $n \in \{1, 6, 24, 96, T\}$ on CASAS hh113
($T = 12{,}226$, $W = 24$, $\theta = 10$, $k = 4$)
and measured total processing time.

\begin{table}[t]
\caption{Processing speed vs.\ scan frequency ($n$) on CASAS hh113.}
\begin{center}
\begin{tabular}{|c|r|r|}
\hline
$n$ & Time (ms) & Speedup vs.\ $n{=}1$ \\\hline
1   & 1,068     & 1.0$\times$ \\
6   &   216     & 4.9$\times$ \\
24  &    69     & 15.5$\times$ \\
96  &    27     & 39.6$\times$ \\
$T$ (batch) & 7.5 & 142.4$\times$ \\\hline
\end{tabular}
\end{center}
\end{table}

$n = W = 24$ achieves a 15.5$\times$ speedup over $n = 1$
while scanning every 24 hours, a practically useful frequency.
The improvement saturates after $n = W$,
suggesting $n = W$ as the recommended default.

\subsection{EX-2: Ablation Study}

We measured the contribution of each speedup technique
on synthetic data ($T = 12{,}226$, $W = 24$, $\theta = 10$, $k = 4$).

\begin{table}[t]
\caption{Ablation: processing time [ms] for four algorithm variants.}
\begin{center}
\begin{tabular}{|l|c|c|r|}
\hline
Variant & Stride & Pruning & Time (ms) \\\hline
Naive          &        &         & TODO \\
Stride-only    & \checkmark &    & TODO \\
Pruning-only   &        & \checkmark & TODO \\
Full (proposed)& \checkmark & \checkmark & TODO \\\hline
\end{tabular}
\end{center}
\end{table}

% TODO: fill in ablation results after running experiments

\subsection{EX-3: Application to CASAS Sensor Data}

We applied the algorithm ($n = 24$, batch mode) to two CASAS datasets.
CASAS records multi-sensor co-activation streams structurally analogous
to autonomous robotic behavioral logs.

\begin{table}[t]
\caption{CASAS dataset statistics and detected dense intervals.}
\begin{center}
\begin{tabular}{|l|r|r|r|}
\hline
Dataset & $T$ & $W$ & Patterns \\\hline
hh113 & 12,226 & 24 & 57 \\
mv001 & 34,759 & 24 & 53 \\\hline
\end{tabular}
\end{center}
\end{table}

% TODO: add representative pattern examples and interpretation
